\documentclass[12pt,dvipdfmx]{beamer}

%% for lstcoq
\definecolor{dkgreen}{rgb}{0,0.6,0}
\definecolor{dkviolet}{rgb}{0.6,0,0.4}
\definecolor{dkblue}{rgb}{0,0,0.6}
\definecolor{ltblue}{rgb}{0.2,0.2,1}

\usepackage{graphicx,listings}
\usepackage{sty/bcprules,sty/jlisting,sty/lstcoq} % sty ディレクトリ以下にあるものは sty/hoge のように sty/ をつける
\usepackage{sty/jtygm,type1cm}                       % フォント関連の警告を消す
\usepackage{macros}              % my macros

%% テーマ
\usetheme{default}               % 全体
\usecolortheme{rose}             % 色
\usefonttheme{structurebold}     % タイトルのフォント
\usefonttheme{professionalfonts} % 数式のフォント用

%% リストの形とか
\useinnertheme{circles}

%% 上に出てくる、目次の前、現在、後の名前を表示する
\useoutertheme[subsection=false]{miniframes}

%% 下の方にでる謎なナビゲーションを消す
\setbeamertemplate{navigation symbols}{}

%% ページ番号
\setbeamertemplate{footline}[frame number]
\setbeamerfont{footline}{size=\footnotesize}
\setbeamercolor{footline}{fg=black,bg=black}

%% 日本語フォントをゴシックに
\renewcommand{\kanjifamilydefault}{\gtdefault}
\renewcommand{\familydefault}{\sfdefault}

%% カスタムカラーテーマ (from beaver.sty)
\definecolor{themecolor}{rgb}{0,0.1,0.4} % ここを変えるとタイトルの色とかかわる

\setbeamercolor{section in toc}{fg=black,bg=white}
\setbeamercolor*{palette primary}{fg=themecolor!60!black,bg=gray!30!white}
\setbeamercolor*{palette secondary}{fg=themecolor!70!black,bg=gray!15!white}
\setbeamercolor*{palette tertiary}{bg=gray!10!white,fg=themecolor!20!black}
\setbeamercolor*{palette quaternary}{fg=themecolor,bg=gray!5!white}
\setbeamercolor{titlelike}{parent=palette primary,fg=themecolor} % タイトルの色
\setbeamercolor{frametitle}{bg=gray!10!white}                    % フレームタイトルの色
\setbeamercolor{frametitle right}{bg=gray!60!white}
\setbeamercolor*{separation line}{}
\setbeamercolor*{fine separation line}{}
\setbeamercolor{local structure}{fg=themecolor!80!black} % enumerate, itemize の色

%% listings
%% Coq
\lstdefinestyle{coqstyle}{
  language=Coq,
  breaklines=true,
  basicstyle=\ttfamily\scriptsize,
  frame=tb,
  framesep=6pt,
  captionpos=b
}
%% Erlang
\lstdefinestyle{erlangstyle}{
  language=Erlang,
  breaklines=true,
  basicstyle=\ttfamily\scriptsize,
  frame=tb,
  framesep=6pt,
  captionpos=b
}
%% the others
\lstset{
  breaklines=true,
  basicstyle=\ttfamily\scriptsize,
  frame=tb,
  framesep=6pt,
  captionpos=b
}

%% タイトルページ
\title{My Beamer Template}
\author{Shohei Yasutake}
\institute{Tokyo Institute of Technology}
\date{\today}

%% 右下にロゴ
%% \logo{\includegraphics[width=1.3cm]{./img/psg.png}}

\begin{document}

\maketitle

\begin{frame}{Outline}
  \tableofcontents
\end{frame}

\section{Introduction}

\slide{Introduction}{}{
  This is the first slide.
  \\~\\
  これは自分の Beamer のテンプレートです。
  \\~\\
  空行は \texttt{\textbackslash\textbackslash\textasciitilde\textbackslash\textbackslash} です
}

\section{Body}

\subsection{hoge}

\slide{Itemize}{}{
  Body Section

  wideitemize - itemize
  \begin{wideitemize} % 間がちょっと広い
  \item Coq
  \item Haskell
  \item ML
    \begin{itemize}
    \item OCaml
    \item SML
    \item SML\#
    \end{itemize}
  \item Scala
  \end{wideitemize}
}

\slide{日本語}{}{
  これは日本語です。

  \begin{align}
    f(x) & = x^2 \\
    g(x) & = 2x
  \end{align}
}

\subsection{block 系}

\slide{Block}{block,exampleblock,alertblock}{
  \begin{block}{block}
    \tt{block}
  \end{block}

  \begin{exampleblock}{exampleblock}
    \tt{exampleblock}

    \coqinline{Theorem modus\_ponens : forall P Q, P -> (P -> Q) -> Q}
  \end{exampleblock}

  \begin{alertblock}{alertlock}
    \tt{alertblock}
  \end{alertblock}
}

\slide{Theorem}{theorem,corollary, fact, lemma}{
  \begin{theorem}[すごい定理]
    For any $a \in A$, $a$ はなんかすごい性質を満たす
  \end{theorem}

  \begin{corollary}[すごい系]
    For any $a \in A$, $a$ はなんかすごい性質を満たす
  \end{corollary}

  \begin{fact}[すごい事実]
    For any $a \in A$, $a$ はなんかすごい性質を満たす
  \end{fact}

  \begin{lemma}[すごい補題]
    For any $a \in A$, $a$ はなんかすごい性質を満たす
  \end{lemma}
}

\slide{Definition}{}{
  \begin{definition}[基本的な定義]
    ほげ
  \end{definition}

  \begin{example}[例っぽい]
    「なになにっぽい」ってなんなんでしょうか
  \end{example}

  \begin{proof}[すごい定理の証明です]
    明らか
  \end{proof}

  \begin{proof}
    すごい定理から明らか
  \end{proof}
}

\slide{Inference rules}{}{
  from bcprules.sty placed at \texttt{./sty/}

  \infrule[Nil]{
  }{
    nil : list A
  }

  \infrule[Cons]{
    a : A
    \andalso
    l : list A
  }{
    a :: l : list A
  }
}

\subsection{Source Codes}

\slide{Source code}{Coq}{
  \lstinputlisting[language=Coq,style=default]{./src/coq.v}
}

\slide{Source code}{Haskell}{
  \lstinputlisting[language=Haskell,style=default]{./src/haskell.hs}
}

\slide{Source code}{Erlang}{
  \lstinputlisting[language=Erlang,style=default,caption=Erlang]{./src/erlang.erl}
}

\slide{Source code}{OCaml}{
  \lstinputlisting[language=OCaml,style=default]{./src/ocaml.ml}
}

\begin{frame}[fragile]{lstlisting}{}

  requires \texttt{fragile} option

  \begin{lstlisting}[style=default,caption={[JS]JavaScript}]
    function() {
      hoge;
      poyo;
      return this;
    }
  \end{lstlisting}
\end{frame}

\subsection{fuga}

\slide{Image}{}{
  \begin{figure}
    \includegraphics[width=5cm]{./img/psg.png}
  \end{figure}

  PSG のロゴ\footnote{\url{http://www.psg.cs.titech.ac.jp}}
}

\section{Conclusion}

\slide{Conclusion}{}{
  This is the last slide.
}


\end{document}
